
\subsubsection{2.1 Reward Model Evaluation}

RewardBench provides a benchmark spanning four categories: Chat, Chat Hard, Safety, and Reasoning. It evaluates whether reward models correctly identify the better response in curated pairs. ArmoRM achieves 89.0\% overall accuracy on this benchmark. MRMBench extends reward model probing to multiple dimensions including instruction following, harmlessness, and factual accuracy. Both benchmarks focus on content quality rather than structural formatting preferences.

\subsubsection{2.2 Preference Data and Annotation}

Reward models learn from human preference data collected through pairwise comparisons. Research on annotation artifacts has documented systematic patterns in human labeling that models can learn to exploit. Cognitive factors such as ease of scanning enumerated text or perceived thoroughness of listed content may influence rater preferences during data collection, though this mechanism was not directly tested in the present study.

\subsubsection{2.3 Reward Model Architectures}

Modern reward models employ diverse architectures. Decoder-only models such as ArmoRM, UltraRM, and Starling-RM process tokens sequentially with causal attention. Encoder-only models such as PairRM (based on DeBERTa) use bidirectional attention and pairwise comparison. These architectural differences may affect how models process structural features.
